\section{Introduction}

The urgency of mitigating global warming necessitates a rapid energy transition to achieve the Sustainable Development Goals (SDGs) \cite{messerli2019global}. Energy-system transformation is central to climate-change mitigation \cite{BOGDANOV2021120467}. This transition is not solely a technological endeavor \cite{CHEN20191564}; it also encompasses critical economic dimensions, among which the weighted average cost of capital (WACC) is a critical metric \cite{BACHNER2019100373}. It represents the average return required by capital providers, balancing the costs of debt and equity financing \cite{BACHNER2019100373}. The environmental consequences of this financial indicator are profound; Ameli et al. report 35\% lower green electricity production in Africa under region-specific WACC assumptions in a cost-optimal 2$^\circ$C pathway; their financing-policy scenario allows Africa to reach net-zero emissions approximately ten years earlier \cite{ameli2021higher}.

This study aligns with the UN Sustainable Development Goals (SDGs), specifically SDG 7 (Affordable and Clean Energy) through research on renewable project finance, SDG 13 (Climate Action) by examining financing conditions relevant to decarbonization, SDG 9 (Industry, Innovation and Infrastructure) through an open, reproducible three-stage machine learning pipeline, and SDG 17 (Partnerships for the Goals) by releasing harmonized multi-regional panels to support global research collaboration.

A pernicious challenge for the energy transition is its inherent vulnerability to shocks from the legacy fossil fuel system it seeks to replace. Global oil, gas, and coal price fluctuations can affect inflation and financing conditions, creating potential exposure for capital-intensive renewable projects \cite{melodia2022energy}. Financing assumptions can change modeled investment choices and transition pathways \cite{egli2019bias}. Existing research has examined the relationship between WACC, interest rates, and unstable market conditions in the energy sector \cite{dobrowolski2022weighted}. However, less attention has been devoted to evaluating its influence on WACC within the context of energy transition projects. Additionally, the varied impacts of fossil fuel price volatility across different energy sectors and regions have not been fully explored.

Therefore, this study focuses on the upcoming decade and seeks to answer the following questions:
\begin{itemize}
    \item How do fossil fuel price fluctuations relate to the WACC of energy transition projects in modeled future scenarios?
    \item How do these price dynamics relate to capital costs across regions, particularly between fuel-importing and fuel-exporting countries, and how should financing strategies for low-carbon projects be shaped?
\end{itemize}

This work contributes an open and reproducible AI framework that translates fossil fuel price signals into actionable climate finance intelligence. By organizing regional, fuel, and project comparisons within an AutoGluon-based forecasting pipeline, the framework provides decision support for analyzing regional differences in financing conditions, energy costs, and decarbonization pathways. Its outputs can inform efforts to reduce household energy burdens, improve environmental quality, and support more effective climate and energy policy across diverse global regions.



\textbf{Open Science Statement:} Code, data-processing instructions, cleaned WACC scenario panels, archived forecast comparisons, and corrected sensitivity diagnostics for this version are available in the organization repository\footnote{\url{https://github.com/GreenComp-ERC/ImpactOfFossilFeulPrice}}. An archived dataset snapshot is hosted on Hugging Face\footnote{\url{https://huggingface.co/datasets/global-nomad-nexus/ImpactOfFossilFuelPrice-WACC}}; the repository records the version and provenance of the diagnostics used in this paper. Editable paper, slides, and poster sources are provided in the repository's \texttt{conference/ACAIT2026/} directory.
